\section*{अध्याय \ref{ch:g4:volcanoes-earthquakes-first-look} --- ज्वालामुखी और भूकंप: पहली झलक}

\begin{solution}{exo:g4:volcanoes-earthquakes-first-look:1}
(क) \omterm{def:g4:volcanoes-earthquakes-first-look:eruption}{लावा}। (ख) \omterm{def:g4:volcanoes-earthquakes-first-look:eruption}{ज्वालामुखीय राख}। (ग) गैस।
\end{solution}

\begin{solution}{exo:g4:volcanoes-earthquakes-first-look:2}
\omterm{def:g4:volcanoes-earthquakes-first-look:earthquake}{भूकंप} ज़मीन का अचानक काँपना है, जो ज़मीन के बहुत नीचे टूटती और खिसकती
\omterm{def:g3:stones-pebbles-sand:rock}{चट्टानों} के कारण होता है। वह सेकंडों तक चलता है (सबसे बड़े \omterm{def:g4:volcanoes-earthquakes-first-look:earthquake}{भूकंपों} में एक-दो
मिनट), घंटों तक नहीं।
\end{solution}

\begin{solution}{exo:g4:volcanoes-earthquakes-first-look:3}
झुको, ढको, पकड़े रहो।
\end{solution}

\begin{solution}{exo:g4:volcanoes-earthquakes-first-look:4}
प्रशांत \omterm{def:g2:water-on-earth:ocean}{महासागर} के चारों ओर।
\end{solution}

\begin{solution}{exo:g4:volcanoes-earthquakes-first-look:5}
(क) सही: ज़्यादातर \omterm{def:g4:volcanoes-earthquakes-first-look:earthquake}{भूकंपों} को केवल यंत्र पकड़ते हैं। (ख) ग़लत: \omterm{def:g4:volcanoes-earthquakes-first-look:eruption}{लावा} ठंडा होकर
कठोर नई \omterm{def:g3:stones-pebbles-sand:rock}{चट्टान} बन जाता है। (ग) ग़लत: ज़मीन काँपते समय भागना ख़तरनाक है (तुम
गिर सकते हो, और इमारतों से काँच और चीज़ें गिरती हैं); झुको, ढको और
पकड़े रहो।
\end{solution}

\begin{solution}{exo:g4:volcanoes-earthquakes-first-look:6}
वह \omterm{def:g4:volcanoes-earthquakes-first-look:volcano}{ज्वालामुखी} के बहुत नीचे से आती है, जहाँ \omterm{def:g3:stones-pebbles-sand:rock}{चट्टान} पिघल गई है। वह \omterm{def:g4:volcanoes-earthquakes-first-look:volcano}{ज्वालामुखी} के
बीच से होकर ऊपर उठती है और चोटी पर बाहर निकलती है।
\end{solution}

\begin{solution}{exo:g4:volcanoes-earthquakes-first-look:7}
(क) चिली (पेरू, इक्वाडोर या मेक्सिको भी सही हैं): उसके \omterm{def:g4:volcanoes-earthquakes-first-look:volcano}{ज्वालामुखी} प्रशांत \omterm{def:g1:landscapes-around-us:water}{तट} के
पास उसकी पूरी लंबाई के साथ फैले हैं। (ख) हवाई, प्रशांत \omterm{def:g2:water-on-earth:ocean}{महासागर} के बीच में,
छायांकित वलय से दूर।
\end{solution}

\begin{solution}{exo:g4:volcanoes-earthquakes-first-look:8}
$141 \div 12$ 12 से थोड़ा कम है (क्योंकि $12 \times 12 = 144$): महीने में लगभग
12 बड़े \omterm{def:g4:volcanoes-earthquakes-first-look:earthquake}{भूकंप}।
\end{solution}

\begin{solution}{exo:g4:volcanoes-earthquakes-first-look:9}
20 का एक-चौथाई 5 है, इसलिए तीन-चौथाई $3 \times 5 = 15$ है: वलय में लगभग 15,
और उसके बाहर लगभग $20 - 15 = 5$।
\end{solution}

\begin{solution}{exo:g4:volcanoes-earthquakes-first-look:10}
उसे अपने हाथों और घुटनों पर झुकना चाहिए था, अपनी मेज़ के नीचे सिर और गर्दन को
ढकना चाहिए था, और काँपना रुकने तक मेज़ को पकड़े रहना चाहिए था। दरवाज़े की चौखट
उसे गिरती या उड़ती चीज़ों से नहीं बचाती, और ज़मीन काँपते समय भागने से वह गिर
सकता है।
\end{solution}

\begin{solution}{exo:g4:volcanoes-earthquakes-first-look:11}
वलय के बाहर: $1214 - 837 = 377$ \omterm{def:g4:volcanoes-earthquakes-first-look:volcano}{ज्वालामुखी}। 1\,200 का एक-तिहाई 400 है, इसलिए
1\,200 का दो-तिहाई 800 है। वलय में 837 हैं, लगभग 1\,200 में से 800 से थोड़े
ज़्यादा: ``लगभग 3 में से 2'' अच्छा वर्णन है।
\end{solution}
